\documentclass[a4paper,12pt]{article}
\usepackage[utf8]{inputenc}
\usepackage[T1]{fontenc}
\usepackage[italian]{babel}
\usepackage{amsmath,amssymb}
\usepackage{booktabs}
\usepackage{geometry}
\geometry{margin=2.5cm}

\begin{document}

\section*{Sviluppo della formulazione QUBO per il Multi-Dimensional Knapsack Problem}

Siano:
\[
x_i \in \{0,1\}, \qquad i=1,\dots,n,
\]
dove \(x_i = 1\) indica che l'oggetto \(i\) viene selezionato, mentre \(x_i = 0\) indica il contrario.

Indichiamo con:
\begin{itemize}
    \item \(p_i\) il profitto dell'oggetto \(i\);
    \item \(w_{id}\) il peso dell'oggetto \(i\) rispetto al vincolo \(d\);
    \item \(C_d\) la capacità associata al vincolo \(d\);
    \item \(\lambda_d\) il coefficiente di penalità associato al vincolo \(d\), con \(d=1,\dots,D\).
\end{itemize}

Per ogni vincolo \(d = 1, \dots, D\):

\begin{itemize}
    \item \(w_{id} \in \mathbb{R}_{\ge 0}\) rappresenta il consumo di capacità dell'oggetto \(i\) rispetto alla dimensione \(d\);
    \item \(C_d \in \mathbb{R}_{>0}\) è la capacità disponibile per il vincolo \(d\);
    \item \(\lambda_d \in \mathbb{R}_{>0}\) è il coefficiente di penalità associato al vincolo \(d\).
\end{itemize}

Dove \(\lambda_d\) è stato fissato come rapporto tra la somma dei profitti degli oggetti e la capacità del vincolo \(d\)-esimo, ossia:
\begin{equation}
\lambda_d = \frac{\sum_{i=1}^{n} p_i}{C_d},
\tag{*}
\end{equation}
per ogni \(d = 1, \dots, D\).

Il problema multidimensional knapsack consiste nel massimizzare il profitto totale rispettando simultaneamente tutti i vincoli di capacità.  
Una corrispondente formulazione penalizzata, espressa come problema di minimizzazione, è la seguente:
\begin{equation}
Q(x)= -\sum_{i=1}^{n} p_i x_i + \sum_{d=1}^{D} \lambda_d \left(\sum_{i=1}^{n} w_{id}x_i - C_d\right)^2.
\end{equation}

Sviluppiamo il termine quadratico relativo a un generico vincolo \(d\):
\begin{align}
\left(\sum_{i=1}^{n} w_{id}x_i - C_d\right)^2
&= \left(\sum_{i=1}^{n} w_{id}x_i\right)^2 - 2C_d\sum_{i=1}^{n} w_{id}x_i + C_d^2 \\
&= \sum_{i=1}^{n} w_{id}^2 x_i^2 + 2\sum_{i<j} w_{id}w_{jd}x_ix_j - 2C_d\sum_{i=1}^{n} w_{id}x_i + C_d^2.
\end{align}

Poiché \(x_i \in \{0,1\}\), vale la proprietà
\[
x_i^2 = x_i.
\]
Inoltre, il termine costante \(C_d^2\) può essere trascurato nella formulazione QUBO, in quanto non influenza l'argomento di minimo. Pertanto:
\begin{equation}
\left(\sum_{i=1}^{n} w_{id}x_i - C_d\right)^2
\sim
\sum_{i=1}^{n} w_{id}^2 x_i + 2\sum_{i<j} w_{id}w_{jd}x_ix_j - 2C_d\sum_{i=1}^{n} w_{id}x_i,
\end{equation}
dove il simbolo \(\sim\) indica equivalenza a meno di una costante additiva.

Sostituendo tale sviluppo nella funzione obiettivo, si ottiene:
\begin{align}
Q(x)
&\sim -\sum_{i=1}^{n} p_i x_i
+ \sum_{d=1}^{D} \lambda_d \left(
\sum_{i=1}^{n} w_{id}^2 x_i
+ 2\sum_{i<j} w_{id}w_{jd}x_ix_j
- 2C_d\sum_{i=1}^{n} w_{id}x_i
\right) \\
&= -\sum_{i=1}^{n} p_i x_i
+ \sum_{d=1}^{D}\sum_{i=1}^{n} \lambda_d w_{id}^2 x_i
+ 2\sum_{d=1}^{D}\sum_{i<j} \lambda_d w_{id}w_{jd}x_ix_j
- 2\sum_{d=1}^{D}\sum_{i=1}^{n} \lambda_d C_d w_{id}x_i.
\end{align}

Raccogliendo i termini lineari in \(x_i\), segue:
\begin{align}
Q(x)
&\sim
\sum_{i=1}^{n}
\left[
-p_i + \sum_{d=1}^{D} \lambda_d w_{id}^2 - 2\sum_{d=1}^{D} \lambda_d C_d w_{id}
\right]x_i
+ 2\sum_{d=1}^{D}\sum_{i<j} \lambda_d w_{id}w_{jd}x_ix_j \\
&=
\sum_{i=1}^{n}
\left[
-p_i + \sum_{d=1}^{D} \lambda_d w_{id}(w_{id}-2C_d)
\right]x_i
+ 2\sum_{d=1}^{D}\sum_{i<j} \lambda_d w_{id}w_{jd}x_ix_j.
\end{align}

Di conseguenza, la formulazione QUBO assume la forma standard:
\begin{equation}
Q(x)=\sum_{i=1}^{n} Q_{ii}x_i + \sum_{i<j} Q_{ij}x_ix_j,
\end{equation}
con coefficienti:
\begin{equation}
Q_{ij}=
\begin{cases}
-p_i + \displaystyle\sum_{d=1}^{D} \lambda_d w_{id}(w_{id}-2C_d), & i=j, \\[1.2ex]
2\displaystyle\sum_{d=1}^{D} \lambda_d w_{id}w_{jd}, & i\neq j.
\end{cases}
\end{equation}

\end{document}